% ============================================================
% 02_related_work.tex — Поврзани истражувања (македонски)
% ============================================================

\section{Методи на Gradient Boosting}

Методите на засилување (boosting) се ансамблови техники во кои голем број слаби предиктори — обично плитки дрва на одлука — се комбинираат секвенцијално за да формираат силен предиктор. Секое наредно дрво се тренира врз резидуалните грешки на претходниот ансамбл, со што постепено се намалува кумулативната загуба. Овој концепт, формализиран во рамката Gradient Boosting Machine на Friedman, прерасна во доминантна парадигма за регресиски и класификациски задачи врз табеларни податоци и константно постигнува врвни резултати на стандардните репери.

\subsection{XGBoost}

Chen и Guestrin \cite{chen2016xgboost} го воведоа XGBoost (eXtreme Gradient Boosting) во 2016 година, систем кој набрзо стана де факто стандард за натпреварите со структурирани податоци. XGBoost воведе неколку клучни иновации во однос на претходните имплементации на boosting: регуларизациски член во функцијата на цел заради спречување на преоптоварување, алгоритам базиран на хистограм за ефикасно пронаоѓање на поделбите, како и поддршка за паралелно и дистрибуирано тренирање. Системот ги гради дрвата ниво по ниво (level-wise), при што ги оценува сите кандидат-поделби на секоја длабочина пред да ги одбере оние со највисок придобивок. Иако овој пристап обезбедува балансирани дрва и предвидливо користење на меморијата, тој може да биде пресметковно расипнички кога мал подмножество листови го носи поголемиот дел од придобивокот.

\subsection{LightGBM}

Ke и соработниците \cite{ke2017lightgbm} го претставија LightGBM (Light Gradient Boosting Machine) во 2017 година, со примарна цел да се надминат ограничувањата на XGBoost врз датасети со голем број карактеристики и тренинг-примероци. Централната иновација е стратегијата за раст на дрвата \textit{leaf-wise}: наместо да ги шири сите листови на една длабочина истовремено (level-wise), алгоритмот секогаш го одбира поединечниот лист со максимален придобивок за поделба, без оглед на длабочината на дрвото. Ова создава подлабоки и асиметрични дрва кои побрзо ја намалуваат загубата по итерација, но можат да преоптоваруваат на мали датасети. Параметрите \texttt{num\_leaves} и \texttt{min\_child\_samples} го ограничуваат растот на дрвата заради управување со овој ризик.

LightGBM дополнително воведе и Gradient-based One-Side Sampling (GOSS), метод кој ги поднабира тренинг-инстанците со мали градиенти (веќе добро научените примери) со намалена стапка, со што се намалуваат пресметковните трошоци без значителна загуба на точност. Exclusive Feature Bundling (EFB) ги групира меѓусебно ексклузивните карактеристики, ефективно намалувајќи ја димензионалноста на просторот на карактеристиките. Споредбата меѓу XGBoost и LightGBM е директно релевантна за оваа теза, бидејќи и двата алгоритма беа евалуирани во кампањата за оптимизација на хиперпараметрите, при што LightGBM беше потврден како супериорниот избор.

\section{Спортска аналитика и Очекувани голови}

Аналитиката во фудбалот помина низ трансформациска еволуција во последното десетлетие, движена од напредокот во технологиите за следење, достапноста на отворени податоци и статистичката методологија. Традиционалните агрегирани статистики — голови, асистенции, проценти на поседување — беа надополнети со напредни метрики кои поточно го отсликуваат квалитетот на индивидуалните и колективните перформанси.

\subsection{Очекувани голови (xG)}

Метриката за очекувани голови (xG) на секој удар му припишува веројатност да резултира со гол врз основа на набљудуваните карактеристики: позицијата на теренот, аголот во однос на голот, растојанието до голот, видот на ударот (со глава наспроти со нога) и природата на асистенцијата. Rathke \cite{rathke2017expected} спроведе формална анализа на xG моделите користејќи податоци од камери за следење, демонстрирајќи дека xG е значително подобар предиктор на идните тимски перформанси отколку вистинскиот број на постигнати голови. Овој наод е особено важен за FPL: тимовите кои конзистентно генерираат висок xG создаваат вистински напаѓачки можности независно дали тие можности случајно завршиле со гол во конкретна натпреварувачка недела.

Во нашиот систем, xG и xGA (Expected Goals Against) на ниво на тим, преземени од Understat, ја формираат основата за карактеристиките на тимскиот форм. Подвижните просеци од по пет натпреварувачки недели даваат стабилен сигнал за напаѓачкиот и одбранбениот квалитет на тимот, кој ги изедначува краткорочните флуктуации.

\subsection{Вреднување на акциите во фудбалот}

Decroos и соработниците \cite{decroos2019vaep} го воведоа VAEP (Valuing Actions by Estimating Probabilities), рамка за припишување вредност на секоја поединечна акција на играчот врз основа на нејзиниот ефект врз веројатноста за постигнување или примање гол. За разлика од xG, кој мери само удари, VAEP вреднува секое додавање, дрибл и позициско движење. Иако VAEP податоците не се директно вградени во оваа теза поради ограничувања во достапноста на датасети, принципот на статистичко вреднување на индивидуалните придонеси на играчите ја поддржува дизајнерската замисла на нашите карактеристики на подвижниот форм на играчот.

\section{Истражувања за предвидување во фантази-спортовите}

Прогнозирањето на перформансите во фантази-спортовите привлече истражувачко внимание уште од почетокот на 2010-тите години. Hunter и соработниците \cite{hunter2012fantasy} беа меѓу пионерите кои го формализираа проблемот за предвидување фантази-поени со помош на машинско учење, демонстрирајќи дека индивидуалните историски статистики поседуваат значителна предиктивна моќ. Клучно, тие исто така покажаа дека составувањето на оптималниот фантази-тим под буџетски и составни ограничувања е NP-тежок комбинаторен проблем и дека лакомите хеуристики често имаат послаби резултати од формалните оптимизациски пристапи. Ова набљудување директно го мотивира користењето на ILP во нашиот систем, наместо лаком избор.

Поновите истражувања во контекстите на Fantasy NFL и Fantasy NBA покажаа дека ансамбловите методи и невронските мрежи можат да ги надминат линеарните модели кога е достапна доволна количина историски податоци. Контекстот на FPL е специфичен на повеќе начини: варијабилни правила за поени по позиција, системот на чипови (кој создава стратегиски зависности преку повеќе натпреварувачки недели), казни за трансфери, како и потребата за управување низ цела сезона, наместо оптимизација за единечен настан. Оваа теза истовремено ги адресира сите овие специфичности.

\section{Целобројно линеарно програмирање за избор на состав}

Dantzig \cite{dantzig1963linear_programming} ги постави теориските темели на линеарното програмирање, метод со огромен практичен дофат. Целобројното линеарно програмирање (Integer Linear Programming, ILP) е негова генерализација во која некои или сите променливи се ограничени на целобројни вредности, што го прави идеално за комбинаторни проблеми на избор. Симплекс-методот и неговите проширувања (branch-and-bound, отсекувачки рамнини) можат да ги решат ILP инстанците до докажлив оптимум — гаранција што хеуристичките методи како генетските алгоритми или лакомото пребарување не можат да ја понудат.

Во контекст на спортскиот тимски менаџмент, ILP е применуван врз избор на постави во дневните фантази-спортови, каде целта е да се максимизираат проектираните поени, под услов да се почитуваат ограничувањата на платата и составот. Ограничувањата на FPL — буџет, позициски состав, лимит по клуб, цени на трансфери — се природно линеарни, што го прави ILP идеална алатка. Нашата имплементација ја користи библиотеката за моделирање PuLP \cite{mitchell2011pulp} во комбинација со решавачот со отворен код CBC (Coin Branch and Cut).

\section{Бајесова оптимизација на хиперпараметри}

Оптимизацијата на хиперпараметрите е критичен чекор во секој пипелајн за машинско учење. Традиционалните методи ја третираат функцијата на перформанси како црна кутија: пребарувањето во мрежа (grid search) ги евалуира сите комбинации од однапред зададена мрежа (што е пресметковно непрактично во високи димензии), додека случајното пребарување (random search) ги извлекува конфигурациите униформно случајно. Bergstra и Bengio (2012) покажаа дека случајното пребарување го надминува пребарувањето во мрежа кога само неколку хиперпараметри се релевантни, но и двата метода ги игнорираат информациите од претходно евалуираните конфигурации.

Snoek и соработниците \cite{snoek2012practical} воведоа практична Бајесова оптимизација за хиперпараметрите на ML моделите, при што Гаусов процес како сурогат се прилагодува врз набљудуваните парови (хиперпараметри, перформанси) и наредната конфигурација за евалуација се одбира преку максимизирање на функцијата на здобивање (Expected Improvement). Ова овозможува конвергенција со суштински помалку евалуации отколку случајното пребарување, што е особено важно кога секоја евалуација бара изведба на полн пипелајн за тренирање и симулација.

Akiba и соработниците \cite{akiba2019optuna} го развија Optuna, рамка за оптимизација на хиперпараметри која го имплементира семплерот Tree-structured Parzen Estimator (TPE). TPE ги моделира $p(x|y < y^*)$ и $p(x|y \geq y^*)$ одделно — соодветно дистрибуциите на конфигурациите кои постигнале добри и слаби перформанси — и преку нив аналитички го пресметува очекуваното подобрување. TPE се справува поефикасно со високодимензионални и условни хиперпараметарски простори отколку Гаусовите процеси, што го прави добро прилагоден на нашиот заеднички оптимизациски проблем за ML моделот и симулаторот.

\section{Големи јазични модели во поддршката на одлучувањето}

Архитектурата Transformer, претставена од Vaswani и соработниците во 2017 година, ја револуционизираше обработката на природен јазик. Моделите изградени врз оваа архитектура во голем размер демонстрираа значителни способности за расудување.

Brown и соработниците \cite{brown2020gpt3} го претставија GPT-3, јазичен модел со 175 милијарди параметри кој покажа учење со малку примери (few-shot learning): способност да извршува нови задачи кога во барањето добива само мал број примери, без ажурирање на тежините. Ова покажа дека доволно големите јазични модели развиваат емергентни способности за расудување кои се генерализираат низ различни домени. Achiam и соработниците \cite{achiam2023gpt4} го проширија ова со GPT-4, демонстрирајќи суштински подобрувања во математичкото расудување, следењето на инструкции и фактуелната точност.

Семејството Gemini \cite{geminiteam2024gemini}, развиено од Google DeepMind, воведе мултимодални способности заедно со силно јазично расудување. Нашиот систем го користи Gemini 2.5 Flash за стратегиски препораки пред рокот (intel\_05), искористувајќи ја неговата способност да синтетизира структурирани статистички табели заедно со неструктуриран текст од прес-конференции и да произведе применливи совети за управување со составот.

Семејството модели Claude, развиено од Anthropic, се користи во слојот за нараториски објаснувања (Стадиум 9). Неговата способност прецизно да следи сложени инструкции за структуриран излез го прави добро прилагоден за генерирање конзистентни објаснувања по натпреварувачка недела од структурираните податоци за одлуки за составот.

\section{Walk-forward валидација за временски редови}

Tashman \cite{tashman2000out} го формализираше тестирањето надвор од тренинг-примерокот за прогностичките модели, утврдувајќи дека перформансите врз тренинг-примерокот системски ги преценуваат перформансите надвор од него, особено за моделите со многу параметри. Правилниот пристап е строга временска поделба: моделот никогаш не смее да набљудува ниту една податочна точка од периодот на евалуација при тренирањето.

Во контекст на FPL ова е особено критично. Користењето на статистики од идните натпреварувачки недели како карактеристики при тренирање модел чија цел е да ги предвиди токму тие натпреварувачки недели претставува директно протекување на податоците, што ги инфлира пријавените метрики на точност. Претходните истражувања во спортската аналитика и предвидувањето во фантази-спортовите понекогаш не успеале да ја почитуваат оваа граница, давајќи резултати кои не се генерализираат во реални услови.

Нашата walk-forward кросвалидација со пет набори (види Глава~3) строго ја спроведува оваа временска граница. Дополнително, сезоната 2025--26 (евалуациската сезона) е целосно повлечена од сите тренинг-податоци, со што симулацијата на сезоната претставува вистински слеп тест.

\section{Идентификувана празнина во литературата}

По прегледот изложен погоре, може да се идентификува јасна празнина во постојната литература. Ниту едно претходно истражување не ги комбинира следниве четири компоненти во единствен автоматизиран систем:

\begin{enumerate}
  \item Предвидување со машинско учење на индивидуалните перформанси на играчот, со walk-forward временска валидација и онлајн ретренирање по секоја натпреварувачка недела;
  \item Оптимизација на составот со целобројно линеарно програмирање со целосна поддршка за чипови, банкирање трансфери и стратегија низ цела сезона;
  \item Пипелајн за разузнавање во реално време базиран на веб-скрапинг и расудување со LLM, кој ги прилагодува предвидувањата пред секој рок;
  \item Автоматизирана симулација на цела сезона со Бајесова оптимизација на хиперпараметрите низ заедничкиот параметарски простор на ML моделот и симулаторот.
\end{enumerate}

Постојните истражувања обично ја адресираат само една или две од овие компоненти изолирано. Оваа теза ја пополнува таа празнина преку интегрирање на сите четири во кохерентен, евалуиран и документиран систем.
